% ── Polský a slovenský model: postsocialistická srovnání ─────────────────────

\paragraph{Uspořádání sociálního dialogu v~Polsku a~na~Slovensku}\label{par:plsk_architektura} vychází ze~společného dědictví centrálně řízeného
hospodářství a~z~obdobné trajektorie transformace odborového hnutí po~roce~1989.
V~typologii této práce proto tvoří společný postkomunistický srovnávací model,
který doplňuje severský, germánský a~liberální model sociálního dialogu.
V~Polsku je~klíčovým aktérem \textit{NSZZ Solidarność} (~$\approx$~600~tis.~členů),
historicky propojené s~politickou pravicí, a~\textit{Ogólnopolskie Porozumienie
Związków Zawodowych} (OPZZ); zaměstnavatelské svazy reprezentuje
\textit{Konfederacja Lewiatan}. Tripartitní platformou je~\textit{Rada Dialogu
Społecznego} (RDS, od~2015 nahrazující dřívější komisi
\textit{Komisja Trójstronna}).~\cite{Hala_SystemySocialnihoDialogu2013}
Slovensko má~obdobně dvojici konfederací \textit{Konfederácia odborových
zväzov} (KOZ~SR) a~\textit{Asociácia zamestnávateľských zväzov a~združení}
(AZZZ); tripartitu zajišťuje \textit{Hospodárska a~sociálna rada SR}.
Vyjednávání v~obou zemích probíhá převážně na~podnikové úrovni; odvětvové
\acs{KSVS} jsou ojedinělé a~jejich rozšíření \emph{erga omnes} je~právně
možné, ale málo využívané.~\cite{Eurofound_KeyDimensionsIR_2023}

\paragraph{Pokrytí \acs{KS} a~odborová organizovanost}\label{par:plsk_pokryti} vykazují v~obou zemích srovnatelný profil
s~\acgen{ČR}: Polsko dosahuje přibližně \SI{14}{\percent} (nejnižší v~EU-27),
Slovensko \SI{35}{\percent} (data \acs{ICTWSS} 2024,
obr.~\ref{fig:eu_pokryti_kv_mapa}).~\cite{OECD_UnionEmployerCoverage_2025}
Odborová organizovanost v~Polsku klesla z~přibližně \SI{19}{\percent} v~roce~2000
na~současných \SI{12}{\percent}; na~Slovensku ze~zhruba \SI{29}{\percent}
na~\SI{10}{\percent} (obr.~\ref{fig:eu_hustota_mapa}). Pokles je~ve~všech třech
postkomunistických zemích strukturálně shodný a~je~dán kombinací privatizace
státních podniků (s~vysokou mírou organizovanosti), ústupu ze~symboliky
předlistopadových odborů (jejíž zachování přechodem ROH/SROH se~stalo
politickým handicapem) a~celkovou individualizací pracovních
vztahů.~\cite{Eurofound_KeyDimensionsIR_2023}~\cite{Hala_SystemySocialnihoDialogu2013}

\paragraph{Klíčové rysy polského a~slovenského modelu}\label{par:plsk_rysy} spočívají v~nízké aktuální vyjednávací síle navzdory odlišným historickým a~právním východiskům. Polsko je~v~regionu výjimečné silnou symbolickou rolí \textit{Solidarność}
v~politickém systému, přesto je~jeho aktuální vyjednávací síla nízká:
\acs{KS} se~uzavírají převážně v~průmyslových konglomerátech a~veřejném
sektoru. Slovensko zákonem č.~2/1991~Sb. (převzatým z~Československa
po~rozdělení) umožňuje rozšíření \acs{KSVS} \emph{erga omnes}; v~praxi je~však
využíváno omezeně. Obě země jsou silně integrovány do~německého a~rakouského
průmyslového řetězce (subdodavatelé pro~automobilový sektor) a~jejich
mzdový vývoj je~odvozen od~exportní konkurenceschopnosti spíše než
od~vnitřní mzdové komprese skrze \acs{KV}.~\cite{Brixova_MICT_CEE_2026}

Charakteristikou všech tří středoevropských ekonomik (\acgen{ČR}, Polska a~Slovenska)
je~tzv.~\emph{past středního příjmu} (\textit{middle-income trap}):
po~rychlé konvergenční fázi 1995--2008 se~tempo dohánění \aclgen{EU} zpomalilo
a~ekonomiky uvízly přibližně na~\SI{75}{\percent}--\SI{85}{\percent} úrovně
\acgen{HDP} průměru \aclgen{EU} v~\acs{PPS}, aniž by~se~jim~dařilo přejít k~modelu
růstu založenému na~vysoké přidané hodnotě a~vysokých
mzdách. \textit{Brixová} ve~svém srovnání
identifikuje slabý sociální dialog jako~jeden ze~strukturálních faktorů
pasti --- nedostatek koordinace mzdového růstu udržuje konkurenceschopnost
prostřednictvím \emph{nákladového kanálu} (nízké mzdy) místo \emph{produktového
kanálu} (inovace, vyšší přidaná hodnota).~\cite{Brixova_MICT_CEE_2026}

Slovensko vykazuje navzdory nízkému pokrytí \acs{KS} jeden z~nejnižších
Giniho koeficientů v~\acloc{EU} ($\approx 0{,}22$), což je~dáno relativně
homogenní mzdovou strukturou postkomunistické ekonomiky a~plochou daňovou
sazbou. Tento „klid v~rovnosti” je~však křehký: slovenský trh práce vykazuje
zároveň jeden z~nejhlubších regionálních rozdílů v~\acs{EU} (Bratislava versus
východní okresy).~\cite{Eurofound_KeyDimensionsIR_2023}

\paragraph{Implikace pro~\acacc{ČR}}\label{par:plsk_implikace} ukazují ve~srovnání s~Polskem a~Slovenskem, že~\ac{ČR} není v~regionu
ostrovem slabého sociálního dialogu, ale~součástí společného postkomunistického
vzorce. Klíčové ponaučení je~negativní: žádná z~obou zemí nedokázala
za~uplynulých 35~let vybudovat funkční mechanismus rozšiřování \acs{KSVS},
přestože v~obou existoval (nebo nadále existuje) zákonný rámec analogický
českému. Pasivní setrvávání na~současném modelu tedy hrozí stejným
zaseknutím v~pasti středního příjmu, jaké zaznamenává \textit{Brixová}
pro~CEE region jako~celek. Příklad slovenské
nízké mzdové nerovnosti zároveň ukazuje, že~Giniho koeficient sám o~sobě
není dostatečným indikátorem zdraví trhu práce --- rovnost při nízké
absolutní úrovni mezd a~bez kolektivního vyjednávání nese odlišná rizika
než rovnost dosažená koordinovaným \acs{KV} (viz dánský model).~\cite{Brixova_MICT_CEE_2026}
